\section{Introduction}

Paraphrasing defeats contamination detection. When Llama-2-13B is trained on rephrased MMLU items, it achieves 85.9\% accuracy while evading standard 13-gram overlap filters. This exposes a fundamental weakness: surface-form matching cannot detect contamination that preserves semantics but varies phrasing. We investigate whether contamination leaves a behavioral signature that survives paraphrasing---and find a validated mechanism but an inadequate metric.

The problem runs deeper than detection evasion. Current methods occupy two extremes. N-gram overlap, established by GPT-3 and refined by GPT-4's 50-character threshold, scales to trillion-token corpora but fails completely on paraphrased contamination. Quiz-based methods like DCQ achieve higher memorization detection by probing model behavior, but provide only binary outputs and require per-item evaluation, limiting scalability. Neither approach offers what practitioners need: a continuous-valued, scalable, paraphrase-resistant contamination metric.

We hypothesize that contamination produces representation-level changes detectable via confidence behavior, regardless of surface form. The intuition is straightforward: a model trained on benchmark items in various phrasings should develop robust semantic representations that generalize uniformly across paraphrase variations. This uniformity should manifest as consistent confidence scores---low variance---across paraphrases of contaminated items. We formalize this as the Semantic Saturation Index (SSI = 1/variance), where high SSI indicates potential contamination.

Our investigation validates the underlying mechanism but exposes the metric's failure. Through controlled experiments on Mistral-7B with known contamination levels (0--50\% of MMLU), we demonstrate: (1) contamination injection creates measurable training exposure, with 31.1\% accuracy gain at 50\% contamination; (2) paraphrase-augmented training produces representation invariance, with Mean Pairwise Similarity (MPS) differing by 0.065 between training conditions (Cohen's d=0.52); (3) representation invariance correlates with confidence uniformity (r=$-$0.517, d=0.58). The causal chain from contamination to behavioral signature is empirically verified.

However, the SSI metric itself fails to discriminate contamination on real data. On actual MMLU inference with Mistral-7B, SSI achieves AUC=0.506---essentially chance-level classification. The 1/variance formulation amplifies noise rather than signal: within-group variance exceeds between-group differences, and SSI distributions overlap completely across contamination levels. The mechanism is validated; the metric formulation is inadequate.

This paper makes three contributions. First, we provide the first empirical demonstration of the causal chain linking contamination to representation invariance to confidence uniformity. Second, we document the failure of inverse-variance metrics for contamination detection, guiding future work away from simple transforms that amplify noise. Third, we identify a significant simulation-reality gap in contamination research, where simulated results showed PASS across mechanism hypotheses but real-data evaluation revealed metric failure.
